% ═══════════════════════════════════════════════════════════════
% LH-02c: Possessivbegleiter (LEHRERHINWEISE)
% Spanisch Klasse 7 - Sequenz 2: Mi familia
% ═══════════════════════════════════════════════════════════════
% Abgeleitet von: LE-02c.tex (Score: 9.2/10)
% ═══════════════════════════════════════════════════════════════

\documentclass[11pt, a4paper]{article}

% ═══════════════════════════════════════════════════════════════
% SW-MODUS TOGGLE
% ═══════════════════════════════════════════════════════════════
\newif\ifswmodus
\swmodusfalse

% ═══════════════════════════════════════════════════════════════
% PAKETE
% ═══════════════════════════════════════════════════════════════

\usepackage[utf8]{inputenc}
\usepackage[T1]{fontenc}
\usepackage[ngerman]{babel}
\usepackage{helvet}
\renewcommand{\familydefault}{\sfdefault}

\usepackage[margin=2cm]{geometry}
\usepackage{enumitem}
\usepackage{tabularx}
\usepackage{booktabs}
\usepackage{longtable}

\usepackage{xcolor}
\usepackage{amssymb}
\usepackage{tcolorbox}
\tcbuselibrary{skins}

\usepackage{fancyhdr}
\usepackage{lastpage}
\usepackage{needspace}

% ═══════════════════════════════════════════════════════════════
% FARBEN
% ═══════════════════════════════════════════════════════════════

\definecolor{spanisch-color}{HTML}{c41e3a}
\definecolor{grau-color}{HTML}{888888}
\definecolor{hellgrau-color}{HTML}{f5f5f5}
\definecolor{basis-color}{HTML}{2a7a4b}
\definecolor{standard-color}{HTML}{2c5aa0}
\definecolor{erweiterung-color}{HTML}{7b1fa2}
\definecolor{loesung-color}{HTML}{c62828}

\definecolor{sw-dark}{HTML}{000000}
\definecolor{sw-gray}{HTML}{666666}
\definecolor{sw-white}{HTML}{ffffff}

\ifswmodus
    \colorlet{spanisch}{sw-dark}
    \colorlet{grau}{sw-gray}
    \colorlet{hellgrau}{sw-white}
    \colorlet{basis}{sw-dark}
    \colorlet{standard}{sw-dark}
    \colorlet{erweiterung}{sw-dark}
    \colorlet{loesung}{sw-dark}
\else
    \colorlet{spanisch}{spanisch-color}
    \colorlet{grau}{grau-color}
    \colorlet{hellgrau}{hellgrau-color}
    \colorlet{basis}{basis-color}
    \colorlet{standard}{standard-color}
    \colorlet{erweiterung}{erweiterung-color}
    \colorlet{loesung}{loesung-color}
\fi

\newcommand{\fachfarbe}{spanisch}

% ═══════════════════════════════════════════════════════════════
% VARIABLEN
% ═══════════════════════════════════════════════════════════════

\newcommand{\thema}{Possessivbegleiter}
\newcommand{\sequenz}{02}
\newcommand{\block}{c}
\newcommand{\fach}{Spanisch}
\newcommand{\klasse}{7}
\newcommand{\dauer}{90 Min. (Doppelstunde)}
\newcommand{\lerngruppengroesse}{24}
\newcommand{\datum}{\today}

% ═══════════════════════════════════════════════════════════════
% KOPF- UND FUSSZEILE
% ═══════════════════════════════════════════════════════════════

\pagestyle{fancy}
\fancyhf{}
\renewcommand{\headrulewidth}{0.4pt}
\renewcommand{\footrulewidth}{0.4pt}

\lhead{\fach{} -- Klasse \klasse}
\chead{\textbf{LH-\sequenz\block{} (Lehrerhinweise)}}
\rhead{\datum}

\lfoot{}
\cfoot{}
\rfoot{Seite \thepage\ von \pageref{LastPage}}

% ═══════════════════════════════════════════════════════════════
% BEFEHLE
% ═══════════════════════════════════════════════════════════════

\newcommand{\B}{\textcolor{basis}{\textbf{[B]}}}
\newcommand{\Std}{\textcolor{standard}{\textbf{[S]}}}
\newcommand{\E}{\textcolor{erweiterung}{\textbf{[E]}}}
\newcommand{\loesung}[1]{\textcolor{loesung}{\textbf{#1}}}
\newcommand{\es}[1]{\textcolor{\fachfarbe}{\textbf{#1}}}

\newtcolorbox{kurzplan}{
    colback=hellgrau,
    colframe=\fachfarbe,
    colbacktitle=white,
    coltitle=\fachfarbe,
    boxrule=1pt,
    arc=0pt,
    fonttitle=\bfseries,
    attach boxed title to top left={yshift=-2mm, xshift=5mm},
    boxed title style={boxrule=0pt, colback=white},
    title=Kurzplan
}

\newtcolorbox{differenzierung}{
    colback=white,
    colframe=grau,
    colbacktitle=white,
    coltitle=grau,
    boxrule=0.5pt,
    arc=0pt,
    fonttitle=\bfseries,
    attach boxed title to top left={yshift=-2mm, xshift=5mm},
    boxed title style={boxrule=0pt, colback=white},
    title=Differenzierung
}

\newtcolorbox{fehlerbox}{
    colback=white,
    colframe=loesung,
    colbacktitle=white,
    coltitle=loesung,
    boxrule=0.5pt,
    arc=0pt,
    fonttitle=\bfseries,
    attach boxed title to top left={yshift=-2mm, xshift=5mm},
    boxed title style={boxrule=0pt, colback=white},
    title=Häufige Fehler
}

\newtcolorbox{materialbox}{
    colback=white,
    colframe=\fachfarbe,
    colbacktitle=white,
    coltitle=\fachfarbe,
    boxrule=0.5pt,
    arc=0pt,
    fonttitle=\bfseries,
    attach boxed title to top left={yshift=-2mm, xshift=5mm},
    boxed title style={boxrule=0pt, colback=white},
    title=Material-Checkliste
}

% ═══════════════════════════════════════════════════════════════
% DOKUMENT
% ═══════════════════════════════════════════════════════════════

\begin{document}

% ═══════════════════════════════════════════════════════════════
% KOPFBEREICH
% ═══════════════════════════════════════════════════════════════

\begin{center}
    {\Large\bfseries\textcolor{\fachfarbe}{Lehrerhinweise: \thema}}\\[0.3em]
    {\normalsize LH-\sequenz\block{} | \dauer}
\end{center}

\vspace{10pt}

% ═══════════════════════════════════════════════════════════════
% 1. KURZPLAN
% ═══════════════════════════════════════════════════════════════

\section*{1. Stundenverlauf (Kurzplan)}

\textbf{1. Stunde (45 Min.) -- Schwerpunkt: mi/tu/su (Singular)}

\begin{tabularx}{\textwidth}{|c|l|X|l|}
\hline
\textbf{Zeit} & \textbf{Phase} & \textbf{Inhalt} & \textbf{Material} \\
\hline
0--5 & Einstieg & Familienfoto: ``Mi familia'' -- ``¿Y tu familia?'' & Bild \\
\hline
5--15 & Input & Possessiv einführen: mi/tu/su + Tabelle & Tafel \\
\hline
15--25 & Übung 1 & Bildkarten: chorisch ``mi madre'', ``tu hermano'' & Karten \\
\hline
25--35 & Übung 2 & AB-02c Aufg. 1 (Lückentext) & AB \\
\hline
35--40 & Sicherung & Lösungen besprechen, Fragen klären & Tafel \\
\hline
40--45 & Überleitung & Plural: mi $\rightarrow$ mis, tu $\rightarrow$ tus, su $\rightarrow$ sus & Tafel \\
\hline
\end{tabularx}

\vspace{1em}

\textbf{2. Stunde (45 Min.) -- Schwerpunkt: Plural + Anwendung}

\begin{tabularx}{\textwidth}{|c|l|X|l|}
\hline
\textbf{Zeit} & \textbf{Phase} & \textbf{Inhalt} & \textbf{Material} \\
\hline
0--5 & Aktivierung & Quiz: ``mi oder mis?'' & -- \\
\hline
5--15 & Übung 3 & AB-02c Aufg. 2 (Singular $\rightarrow$ Plural) & AB \\
\hline
15--25 & Anwendung & AB-02c Aufg. 3 (Familie beschreiben) & AB \\
\hline
25--35 & Dialog & PA: ``¿Cómo se llama tu hermano?'' & AB \\
\hline
35--40 & Sicherung & Merkkasten gemeinsam ausfüllen & AB \\
\hline
40--45 & Abschluss & Zusammenfassung, HA & -- \\
\hline
\end{tabularx}

% ═══════════════════════════════════════════════════════════════
% 2. DIFFERENZIERUNG
% ═══════════════════════════════════════════════════════════════

\section*{2. Differenzierung}

\begin{differenzierung}
\textbf{Legende:} \B{} Basis (BR) | \Std{} Standard (MR) | \E{} Erweiterung (GY)

\vspace{0.5em}

\textbf{WICHTIG:} Diese Marker sind NUR für die Lehrkraft. Auf Schülermaterial erscheinen KEINE Niveaumarkierungen!
\end{differenzierung}

\vspace{1em}

\begin{tabularx}{\textwidth}{|l|X|X|X|}
\hline
& \B{} \textbf{Basis} & \Std{} \textbf{Standard} & \E{} \textbf{Erweiterung} \\
\hline
\textbf{Aufgabe 1} & Possessiv aus Liste wählen & Einsetzen selbstständig & + eigene Beispiele \\
\hline
\textbf{Aufgabe 2} & Tabelle als Hilfe & Transformation selbstständig & + Erklärung der Regel \\
\hline
\textbf{Aufgabe 3} & Mit Satzbausteinen & Mit Redemitteln & Freie Beschreibung \\
\hline
\textbf{Aufgabe 4} & Dialog mit Vorlage & Dialog mit Stichworten & Freier Dialog \\
\hline
\end{tabularx}

\vspace{1em}

\textbf{Konkrete Hinweise:}
\begin{itemize}
    \item \B{} SuS mit LRS: Zeitzugabe (+10 Min.), Possessiv-Übersicht als Hilfe
    \item \B{} DaZ-SuS: Possessiv-Konzept aus L1 aktivieren (sofern vorhanden)
    \item \E{} Leistungsstarke: \textit{nuestro/a} (unser) einführen, Regel selbst formulieren
\end{itemize}

% ═══════════════════════════════════════════════════════════════
% 3. HÄUFIGE FEHLER
% ═══════════════════════════════════════════════════════════════

\section*{3. Häufige Fehler}

\begin{fehlerbox}
\begin{tabularx}{\textwidth}{|l|X|X|}
\hline
\textbf{Fehler} & \textbf{Beispiel} & \textbf{Reaktion} \\
\hline
mi/mis verwechselt & *\es{mi hermanos} statt \es{mis hermanos} & Regel: Singular = mi, Plural = mis; Drill \\
\hline
Genusfehler (*mia) & *\es{mia madre} statt \es{mi madre} & Im Spanisch KEIN Genus bei mi/tu/su! \\
\hline
su-Mehrdeutigkeit & ``su'' = wessen? & Kontext klären: su madre (seine/ihre Mutter) \\
\hline
Transfer aus Deutsch & *\es{meine madre} & Code-Switching vermeiden, spanische Form einüben \\
\hline
tu vs. tú & *\es{Tú hermano} statt \es{Tu hermano} & tu = Possessiv (unbetont), tú = Pronomen (betont) \\
\hline
\end{tabularx}
\end{fehlerbox}

% ═══════════════════════════════════════════════════════════════
% 4. MATERIAL-CHECKLISTE
% ═══════════════════════════════════════════════════════════════

\section*{4. Material-Checkliste}

\begin{materialbox}
\begin{itemize}[label=$\square$]
    \item AB-\sequenz\block.pdf (\lerngruppengroesse{} Kopien)
    \item ML-\sequenz\block.pdf (1× für Lehrkraft)
    \item Lehrwerk: Qué pasa 1, S. 64-66
    \item Bildkarten: Familienmitglieder mit Possessiv
    \item Familienfoto (echtes Foto oder Beispielbild)
    \item Tafelmarker / Kreide
    \item Optional: LP-\sequenz\block.html (Tablets/PCs)
    \item Optional: Stammbaum-Vorlage zum Ausfüllen
\end{itemize}
\end{materialbox}

% ═══════════════════════════════════════════════════════════════
% 5. LÖSUNGEN
% ═══════════════════════════════════════════════════════════════

\section*{5. Lösungen AB-\sequenz\block}

\textbf{Merkkasten: Possessivbegleiter}

\begin{tabular}{llll}
& Deutsch & Singular & Plural \\
yo: & mein/e & \loesung{mi} & \loesung{mis} \\
tú: & dein/e & \loesung{tu} & \loesung{tus} \\
él/ella: & sein/e, ihr/e & \loesung{su} & \loesung{sus} \\
\end{tabular}

\vspace{1em}

\textbf{Aufgabe 1: mi/tu/su einsetzen} (4 Punkte)

\begin{tabular}{ll}
a) & \loesung{Mi} madre se llama Carmen. \\
b) & ¿Cómo se llama \loesung{tu} padre? \\
c) & \loesung{Su} hermano tiene 15 años. \\
d) & ¿Cuántos años tiene \loesung{tu} abuela? \\
\end{tabular}

\vspace{1em}

\textbf{Aufgabe 2: Singular $\rightarrow$ Plural} (4 Punkte)

\begin{tabular}{ll}
a) & mi hermano $\rightarrow$ \loesung{mis hermanos} \\
b) & tu abuelo $\rightarrow$ \loesung{tus abuelos} \\
c) & su prima $\rightarrow$ \loesung{sus primas} \\
d) & mi tío $\rightarrow$ \loesung{mis tíos} \\
\end{tabular}

\vspace{1em}

\textbf{Aufgabe 3: Familie beschreiben} (6 Punkte)

Beispielantworten:
\begin{itemize}
    \item \loesung{Mi madre se llama Carmen. Tiene 42 años.}
    \item \loesung{Mi padre se llama Pedro. Es profesor.}
    \item \loesung{Mis hermanos se llaman Luis y Ana.}
\end{itemize}

\textit{Bewertung: Je 2P pro korrektem Satz (Possessiv + Verb + Info)}

\vspace{1em}

\textbf{Aufgabe 4: Dialog} (6 Punkte)

Mögliche Antworten:
\begin{itemize}
    \item A: ¿Cómo se llama \loesung{tu} hermano?
    \item B: \loesung{Mi} hermano se llama \loesung{Pablo}.
    \item A: ¿Cuántos años tiene?
    \item B: Tiene \loesung{doce} años.
    \item A: ¿Y \loesung{tu} madre?
    \item B: \loesung{Mi} madre se llama \loesung{Ana}.
\end{itemize}

\textit{Bewertung: Fragen korrekt (2P), Antworten mit Possessiv (2P), Namen/Alter (2P)}

% ═══════════════════════════════════════════════════════════════
% 6. HAUSAUFGABE
% ═══════════════════════════════════════════════════════════════

\section*{6. Hausaufgabe}

\begin{itemize}
    \item Lehrwerk S. 65, Übung 5 (Possessivbegleiter einsetzen)
    \item Vokabeln lernen: Possessivbegleiter + 5 Familienwörter
    \item Optional: Lernpfad LP-\sequenz\block.html (interaktive Übungen)
    \item Kreativ: Stammbaum der eigenen Familie zeichnen und beschriften
\end{itemize}

% ═══════════════════════════════════════════════════════════════
% 7. REFLEXIONSNOTIZEN
% ═══════════════════════════════════════════════════════════════

\section*{7. Reflexionsnotizen (nach Durchführung)}

\begin{tabularx}{\textwidth}{|l|X|}
\hline
\textbf{Was lief gut?} & \\[2em]
\hline
\textbf{Was lief nicht?} & \\[2em]
\hline
\textbf{Zeitmanagement} & \\[2em]
\hline
\textbf{Für nächstes Mal} & \\[2em]
\hline
\end{tabularx}

\end{document}
